\documentclass[10pt,a4paper]{article}
\usepackage[utf8]{inputenc}
\usepackage{portuges}
\usepackage{geometry}
\usepackage{graphicx}
\usepackage{booktabs}
\usepackage{longtable}
\usepackage{colortbl}
\usepackage{xcolor}
\usepackage{hyperref}

\geometry{a4paper, margin=2.5cm}

\definecolor{verde}{RGB}{40,167,69}
\definecolor{azul}{RGB}{0,123,255}
\definecolor{cinza}{RGB}{108,117,125}

\title{\textbf{SOC/CISO Dashboard - Análise Comparativa}\\ \large Método Fixed vs Método Adaptativo}
\author{Análise de Logs do Sistema}
\date{Execuções: 17 e 20 de Maio de 2026}

\begin{document}

\maketitle

\begin{abstract}
Este documento apresenta uma análise comparativa entre duas abordagens de gerenciamento de janelas de análise no Agent 1 - Cognitive Security Engine: o método \textbf{Fixed} (janelas de tamanho fixo) e o método \textbf{Adaptativo}. A comparação é baseada em logs reais do sistema SOC/CISO Dashboard v4.2.0, demonstrando as diferenças de comportamento, eficiência e capacidade de detecção entre as duas estratégias.
\end{abstract}

\section{Visão Geral dos Métodos}

\subsection{Método Fixed (Janela Fixa)}
No método Fixed, o sistema utiliza janelas de análise de tamanho constante, sem ajustes dinâmicos baseados no fluxo de alertas.

\subsection{Método Adaptativo}
No método Adaptativo, o WindowManager ajusta dinamicamente o tamanho e o comportamento das janelas de análise conforme o padrão dos alertas recebidos.

\begin{table}[h]
\centering
\caption{Configurações Gerais do Sistema}
\begin{tabular}{@{}lcc@{}}
\toprule
\textbf{Parâmetro} & \textbf{Fixed} & \textbf{Adaptativo} \\
\midrule
Janela CISO base & 15 minutos & 15 minutos \\
Threshold confirmação & 0.85 & 0.85 \\
Máx. hipóteses & 3 & 3 \\
Modo operação & Streaming & Streaming \\
Dataset & 37 alertas & 37 alertas \\
\bottomrule
\end{tabular}
\end{table}

\section{Análise Comparativa Detalhada}

\subsection{Inicialização do WindowManager}

A principal diferença na inicialização está explícita nos logs:

\begin{table}[h]
\centering
\caption{Configuração do WindowManager nos Logs}
\begin{tabular}{@{}p{0.3\textwidth}p{0.6\textwidth}@{}}
\toprule
\textbf{Método} & \textbf{Log de Inicialização} \\
\midrule
\cellcolor{azul!10} Fixed & \texttt{[WindowManager] mode=fixed} (implícito/não evidenciado) \\
\cellcolor{verde!10} Adaptativo & \texttt{[WindowManager] mode=adaptive} (explicitamente registrado) \\
\bottomrule
\end{tabular}
\end{table}

\subsection{Comportamento das Janelas de Análise}

\begin{longtable}{@{}p{0.2\textwidth}p{0.35\textwidth}p{0.35\textwidth}@{}}
\caption{Comparação de Comportamento das Janelas}\\
\toprule
\textbf{Característica} & \textbf{Método Fixed} & \textbf{Método Adaptativo} \\
\midrule
\endfirsthead

\multicolumn{3}{c}{{\tablename\ \thetable{} -- continuação}} \\
\toprule
\textbf{Característica} & \textbf{Método Fixed} & \textbf{Método Adaptativo} \\
\midrule
\endhead

\bottomrule
\multicolumn{3}{r}{{Continua na próxima página}} \\
\endfoot

\bottomrule
\endlastfoot

Padrão de IDs & Sequenciais e previsíveis & Aparentemente aleatórios (ex: \texttt{4c0d4fb5}) \\
Capacidade adaptativa & Baixa & Alta \\
Resposta a padrões repetidos & Pode gerar múltiplas confirmações & Adapta após primeira detecção \\
Fases detectadas & EXFILTRATION → IMPACT & EXFILTRATION → IMPACT \\
Phase Score & 15 & 15 \\
\end{longtable}

\subsection{Processamento de Alertas}

\begin{table}[h]
\centering
\caption{Comparação de Processamento de Alertas}
\begin{tabular}{@{}lcc@{}}
\toprule
\textbf{Métrica} & \textbf{Fixed} & \textbf{Adaptativo} \\
\midrule
Total de alertas & 37 & 37 \\
Tempo total de processamento & $\approx$7 minutos & $\approx$7 minutos \\
Intervalo médio entre alertas & $\approx$11-12 segundos & $\approx$11-12 segundos \\
Alertas até primeira confirmação & 13 & 13 \\
Tempo até primeira confirmação & $\approx$2min23s & $\approx$2min23s \\
\bottomrule
\end{tabular}
\end{table}

\subsection{Hipóteses Confirmadas}

\subsubsection{Método Fixed}
O método Fixed apresentou \textbf{múltiplas confirmações} de diferentes tipos de hipóteses:
\begin{itemize}
\item Advanced Persistent Threats (APT) - 85\%
\item Lateral Movement with Malware - 85\% (múltiplas ocorrências)
\item Ransomware \& Digital Extortion - 85\%
\item APT Data Theft - 85\%
\item Malware Deployment - 85\%
\item Targeted Intrusion Campaign - 85\%
\item Major Breach with Impact - 85\%
\end{itemize}

\subsubsection{Método Adaptativo}
O método Adaptativo apresentou \textbf{apenas uma confirmação}:
\begin{itemize}
\item Major Breach with Impact - 85\% (auto)
\end{itemize}

\begin{quote}
\color{blue}
\textbf{Observação Crítica:} Apesar de ambas as execuções terem identificado duas janelas com o mesmo padrão (EXFILTRATION → IMPACT, 3 alertas, phase\_score=15), o método Adaptativo confirmou apenas a primeira janela, demonstrando capacidade de adaptação e evitando falsos positivos.
\end{quote}

\section{Agent 2 - Graph Reporter}

O comportamento do Agent 2 foi consistente em ambas as execuções, com pequenas variações nos tokens do LLM:

\begin{table}[h]
\centering
\caption{Comparação do Agent 2}
\begin{tabular}{@{}lcc@{}}
\toprule
\textbf{Métrica} & \textbf{Fixed} & \textbf{Adaptativo} \\
\midrule
Provider & Anthropic & Anthropic \\
Threshold & 0.85 & 0.85 \\
Enriquecimento (assets/CVEs/IoCs) & 3/0/2 & 3/0/2 \\
LLM Tokens (prompt/completion) & 465/776 & 465/843 \\
Tempo de análise LLM & $\approx$8s & $\approx$8s \\
Total processamento & $\approx$10s & $\approx$10s \\
\bottomrule
\end{tabular}
\end{table}

\section{Análise de Performance do Orquestrador}

O orquestrador (\_\_main\_\_) apresentou comportamento idêntico em ambas as execuções:

\begin{table}[h]
\centering
\caption{Performance do Orquestrador}
\begin{tabular}{@{}lcc@{}}
\toprule
\textbf{Componente} & \textbf{Tempo de bootstrap} & \textbf{Porta} \\
\midrule
SSE Server & 2 segundos & 5001 \\
Texer & 2 segundos & 5002 \\
Backend FastAPI & 9 segundos & 8000 \\
Agent 1 callback & Imediato & 8002 \\
Tempo total de orquestração & 42 segundos & - \\
\bottomrule
\end{tabular}
\end{table}

\section{Vantagens e Desvantagens}

\subsection{Método Fixed}

\begin{tabular}{@{}p{0.45\textwidth}p{0.45\textwidth}@{}}
\toprule
\textbf{Vantagens} & \textbf{Desvantagens} \\
\midrule
Comportamento previsível & Pode gerar múltiplas confirmações \\
Fácil de depurar & Menos eficiente em padrões repetidos \\
Janelas consistentes & Não se adapta ao contexto \\
\bottomrule
\end{tabular}

\subsection{Método Adaptativo}

\begin{tabular}{@{}p{0.45\textwidth}p{0.45\textwidth}@{}}
\toprule
\textbf{Vantagens} & \textbf{Desvantagens} \\
\midrule
Evita falsos positivos & Comportamento menos previsível \\
Ajusta-se dinamicamente ao fluxo & Mais complexo de depurar \\
Mais eficiente para ataques em fases & Requer lógica adaptativa robusta \\
\bottomrule
\end{tabular}

\section{Conclusão}

A análise comparativa demonstra que ambos os métodos são capazes de processar corretamente os 37 alertas do dataset e detectar o incidente crítico "Major Breach with Impact". No entanto:

\begin{itemize}
\item \textbf{Método Fixed}: Detectou e reportou \textbf{10+ hipóteses diferentes}, gerando múltiplos relatórios.
\item \textbf{Método Adaptativo}: Detectou e reportou \textbf{apenas a hipótese principal}, evitando redundâncias.
\end{itemize}

\vspace{0.5cm}

\boxed{
\begin{minipage}{0.9\textwidth}
\centering
\textbf{\Large Recomendação:}
\vspace{0.3cm}

O \textbf{Método Adaptativo} é superior para ambientes de produção por sua capacidade de reduzir \textbf{falsos positivos} e \textbf{ruído}, focando os recursos da equipe de segurança nos incidentes mais relevantes.
\end{minipage}
}

\section*{Apêndice: Evidências dos Logs}

\subsection*{Log do Modo Adaptativo}
\begin{verbatim}
2026-05-20 14:09:52 [INFO] [agent1.window_manager] 
[WindowManager] mode=adaptive
\end{verbatim}

\subsection*{Log da Confirmação (Modo Adaptativo)}
\begin{verbatim}
2026-05-20 14:12:20 [INFO] [agent1.hypothesis_graph] 
*** AUTO-CONFIRMED: Major Breach with Impact (score=0.850, evidence=1) ***
\end{verbatim}

\subsection*{Log da Segunda Janela (sem confirmação)}
\begin{verbatim}
2026-05-20 14:12:35 [INFO] [agent1.main] 
Window evidence | window_id=8570648a | phases=['EXFILTRATION', 'IMPACT'] 
| phase_score=15 | alert_count=3
\end{verbatim}

\end{document}